% !TeX TS-program = xelatex
% ============================================================================
% Suiyang Guo (YoungYang) - Resume (English)
% Based on adongwanai/LLM-Resume-Template (CC BY 4.0)
% ============================================================================
\documentclass{resume-photo}

% Fonts (FandolSong ships with TeX Live; used only for any CJK glyphs)
\usepackage{xeCJK}
\setCJKmainfont{FandolSong-Regular.otf}[BoldFont=FandolSong-Bold.otf, ItalicFont=FandolKai-Regular.otf]
\usepackage{fontspec}

\usepackage{qrcode}

\usepackage{enumitem}
\usepackage{needspace}
\setlist[itemize]{itemsep=1pt, topsep=1pt, parsep=1pt}
\linespread{1.06}

% 重排条目：第一行「标题(粗) + 右对齐日期」，副标题另起一行 —— 日期严格右对齐、永不换行
\RenewDocumentCommand{\ResumeItem}{omoo}{%
  \Needspace{4\baselineskip}%
  \par\vspace{2pt}\noindent
  {\heiti\zihao{-4} #2}%
  \IfValueT{#4}{\hfill{\normalfont\zihao{5} #4}}%
  \par\nopagebreak
  \IfValueT{#3}{{\zihao{5} #3}\par}%
}
\ctexset{section/beforeskip=0.5em, section/afterskip=0.4em}

% ====== Personal info ======
\ResumeName{Suiyang Guo}
\ResumePhoto{photo.jpg}

\begin{document}

\ResumeContacts{
  +86 158-9334-6463,%
  \ResumeUrl{mailto:1327217178@qq.com}{1327217178@qq.com},%
  \textnormal{Northwestern Polytechnical Univ. | MSc, Control Eng. | 2024.09—2027.06}%
}

\ResumeTitle

% ====================== Summary ======================
\noindent Control Engineering MSc candidate at Northwestern Polytechnical University (Class of 2027), focused on dual-arm deformable-object manipulation. Hands-on experience across the VLA lifecycle, from data curation, SFT/offline post-training and evaluation audits to ROS2 real-robot deployment. Optimized $\pi$0.5, X-VLA and FastWAM systems, reducing step latency from 256.5 ms to 32.05 ms (8$\times$) and idle end-effector jitter by 81\%. Skilled at isolating model, data and robot-stack failures through cross-date evaluation, visual ablations and deployment safety gates.

% ====================== Education ======================
\section{Education}
\ResumeItem{Northwestern Polytechnical University (NPU)}
[\textnormal{MSc, Control Engineering | Robotics \& AI | GPA: 89.56/100}]
[2024.09—2027.06]

\ResumeItem{Zhengzhou University}
[\textnormal{BEng, Automation (Elite Class) | Robotics \& AI | GPA: 3.54/4.0}]
[2020.09—2024.06]

% ====================== Skills ======================
\section{Technical Skills}
\begin{itemize}
  \item \textbf{VLA / Robot Learning}: $\pi$0.5, X-VLA, flow matching, action chunking, SFT, imitation learning, AWBC / RECAP, dataset curation, cross-domain evaluation, FastWAM world-action models
  \item \textbf{Model \& Inference Systems}: PyTorch (DDP), JAX/Flax (FSDP), DINOv3 / SigLIP; Triton kernels, CUDA Graphs, torch.compile, RTC and end-to-end latency profiling
  \item \textbf{Programming}: C/C++ (proficient), Python (proficient), MATLAB (fluent); Linux driver/application development
  \item \textbf{Robotics}: ROS/ROS2, dual-arm manipulation and bring-up, IK, gripper/action mapping, MuJoCo, navigation, SLAM, motion planning and point clouds
  \item \textbf{Embedded \& Hardware}: Yocto / i.MX6ULL embedded Linux, STM32F4, FreeRTOS, motor control and PCB design
  \item \textbf{Engineering}: Git, Linux, WebSocket, PyQt5 host applications and Autodesk Inventor
\end{itemize}

% ====================== Experience ======================
\section{Experience}

\ResumeItem{IDEA Research: Visincept}
[\textnormal{VLA \& World Model Intern, Embodied AI Team}]
[2026.04—Present]

\begin{itemize}
  \item Built a dual-Piper cloth-manipulation pipeline on openpi/$\pi$0.5 spanning data curation, SFT, offline evaluation and ROS2 deployment; reduced native-validation MAE@50 from 0.0337 to 0.0079 and exposed a 3.2$\times$ cross-date generalization degradation through a separate audit split
  \item Led VLA inference/control optimization using RTC chunking, torch.compile, CUDA Graphs and Triton autotuning; reduced step latency from 256.5 ms to 32.05 ms (8$\times$), while minimum-jerk/EMA smoothing cut idle end-effector jitter from 3.33 mm to 0.63 mm ($-$81\%)
  \item Brought up X-VLA on the dual-arm platform; designed visual-ablation gates that traced a silent video-path fallback to black frames, raising the vision-to-proprioception sensitivity ratio from 0 to 5.17 after the fix; firmware-level IK reduced the control cycle from 1.4 s to 0.33 s
  \item Implemented an AWBC/RECAP-advantage offline post-training pipeline and adapted RECAP/RLT to a robotic nail-painting task; reproduced FastWAM with WebSocket/ROS2, continuous-gripper support and preflight safety gates; built CRAVE milestone/value signals with 0.865 stage-progress correlation
\end{itemize}

\ResumeItem{Zhongqin Herui (Shaanxi) Technology Co., Ltd.}
[\textnormal{Navigation \& Mapping Algorithm Engineer (Part-time)}]
[2025.11—2026.01]

\begin{itemize}
  \item Developed 3D mapping and navigation for an orchard autonomous inspection vehicle
  \item Owned equipment/sensor selection, calibration, and sensor-data processing
  \item Built 3D maps from raw LiDAR, then occupancy grids for vehicle navigation
\end{itemize}

\ResumeItem{Huawei Digital Energy Bootcamp}
[\textnormal{Trainee (Training)}]
[2025.07]

\begin{itemize}
  \item Studied digital-energy cloud microservice (DDD) architecture and AI applications/engineering in energy services
  \item Practiced prompt engineering and DevOps deployment workflows
\end{itemize}

% ====================== Projects ======================
\section{Projects}

\ResumeItem{Series-Parallel Hybrid Wheeled Humanoid Robot}
[\textnormal{National Key Research Project | Key Provincial Lab — team member}]
[2024.09—Present]

\begin{itemize}
  \item Built a wheeled humanoid comprising an AGV base, parallel mechanism, dual 6-DoF arms and a 2-DoF head; created a MuJoCo simulation/control-integration environment and a ROS master controller for mapping, planning, obstacle avoidance, point-cloud processing and IK
  \item Developed in-house dual-arm and chassis motor-control systems and integrated STM32F4 firmware with actuators; built a Yocto/i.MX6ULL embedded-Linux image and ROS multi-machine networking, doubling the arm control frequency
\end{itemize}

\ResumeItem{AI-LLM Voice-Interaction Robot}
[\textnormal{Outstanding Undergraduate Thesis | Solo developer}]
[2023.12—2024.09]

\begin{itemize}
  \item Independently built a Raspberry Pi/PyQt5/ROS voice robot integrating audio capture, ASR, LLM dialogue, speech synthesis, digital-human lip sync and smart-home control; received 2 software copyrights
\end{itemize}

\ResumeItem{TCM Therapy Device Series (3 generations)}
[\textnormal{"Challenge Cup" Provincial 2nd Prize | Core R\&D member}]
[2021.03—2023.09]

\begin{itemize}
  \item \textbf{Products}: smart gourd-moxibustion (gen 1) → warm-needle therapy (gen 2) → home mobile moxibustion (gen 3)
  \item \textbf{Stack}: STM32, FreeRTOS, fuzzy PID, sensor fusion, Bluetooth
  \item \textbf{Work}: led full-stack R\&D across three generations (mechanical → circuit → embedded → host software); STM32 firmware (FreeRTOS, fuzzy-PID, sensor integration); hardware circuits (Bluetooth, servo drive, smoke handling); PCB design; multi-sensor (temperature/image/pressure) integration; servo \& stepper motor drivers
\end{itemize}

% ====================== Awards ======================
\section{Awards \& Honors}
\begin{itemize}
  \item "Challenge Cup" Henan Collegiate Academic Sci-Tech Works Competition (16th) \hfill Provincial 2nd Prize (2023)
  \item 2nd Collegiate Electrical \& Electronic Engineer Innovation Contest \hfill Central-China 1st Prize (2023)
  \item 16th "Siemens Cup" Intelligent Manufacturing Challenge \hfill Western Region 2nd Prize (2022)
  \item Lanqiao Cup — Embedded Design \& Development \hfill Provincial 3rd Prize (2022)
  \item HIT Zhengzhou Research Institute Training Camp \hfill Outstanding Camper (2022)
  \item Zhengzhou Univ. 2nd-Class Scholarship \& Merit Student \hfill 2022—2023, 2024—2025
\end{itemize}

% ====================== Patents ======================
\section{Patents \& Certificates}
\begin{itemize}
  \item 4 invention patents, 2 utility-model patents, 2 software copyrights
  \item CET-6 (College English Test Band 6): 473
\end{itemize}

\vfill
\begin{center}
  \qrcode[height=2cm]{https://youngyang-gthb.github.io/}\\[4pt]
  {\footnotesize Scan for online portfolio (projects \& details): \ResumeUrl{https://youngyang-gthb.github.io/}{youngyang-gthb.github.io}}
\end{center}

\end{document}
